% Alpha Decay: Measuring Momentum's Shrinking Half-Life — research-note PDF.
%
% PROSE AND CODE ARE VERBATIM from the website article:
%   site/src/content/articles/alpha-decay-momentum.tsx
%   site/src/content/articles/alpha-decay-code.ts
% Metadata from site/src/lib/articles.ts; project-card fallbacks from
% lib/topics.ts (this article does not override them).
%
% NO CHARTS, deliberately. The article's own header records that this piece is
% unexecuted: the notebook ships without stored outputs, there is no data
% module, and the page shows prose and code only. Every other note draws its
% figures from real computed results; there are none here, so plotting anything
% would mean inventing them.
%
% Layout follows the reference: section heads full width, two-column body,
% code at the full measure.
\documentclass{pyportfolios-note}

\noteshorttitle{Alpha Decay}

\begin{document}
\thispagestyle{notecover}

\notemasthead{Quantitative Insights}{Research Note}{Q3 2026}{}

\notetitle{Alpha Decay: Measuring\\Momentum's Shrinking Half-Life}

\notedeck{The data shows that sector momentum's predictive power fades within
weeks — and has weakened decade over decade.}

\textbf{The finding, upfront.} Using 20 years of daily data on the 11 SPDR
sector ETFs, we measure how long a cross-sectional momentum signal keeps
predicting returns. Two results:

\notecard
  {Algorithmic Trading}
  {Alphalens · Pandas · statsmodels}
  {11 SPDR sector ETFs}
  {Jan 2005 – Dec 2024}
  {Deciding how fast to trade a signal before its edge evaporates — sets
   turnover \& capacity}

\begin{multicols}{2}
\begin{notebulletsnum}
\item \textbf{Horizon decay:} the signal's information coefficient (IC) is
      strongest for the first days after formation and decays roughly
      exponentially — we estimate its half-life directly from the IC curve.
\item \textbf{Calendar decay:} comparing 2005–2014 against 2015–2024, the same
      signal is materially weaker in the recent decade — consistent with the
      \textit{alpha erosion} documented after factor publication (McLean \&
      Pontiff 2016).
\end{notebulletsnum}

For a practitioner both numbers bind: the first sets your \textbf{turnover},
the second your \textbf{expectations}.
\end{multicols}

\begin{notecode}
import numpy as np
import pandas as pd
import matplotlib.pyplot as plt
from scipy import stats
from scipy.optimize import curve_fit

np.random.seed(42)
plt.rcParams["figure.dpi"] = 110
\end{notecode}

% ======================================================== 1 Data & signal ====
\noteneedroom{150bp}
\notesection{1.\notenumsep Data \& signal}

\begin{multicols}{2}
\textbf{Universe:} the 11 GICS sector SPDRs — a small, clean cross-section with
20+ years of history and no survivorship issues.

\textbf{Signal:} classic 12-1 momentum — trailing 252-day return, skipping the
most recent 21 days (to avoid short-term reversal) — recomputed daily,
expressed as cross-sectional ranks.

\textbf{Evaluation:} rank IC — the Spearman correlation between today's signal
ranks and \textit{forward} returns over horizons from 1 to 126 trading days.

\textit{(Alphalens automates exactly this pipeline for larger universes; with 11
assets the direct computation keeps every step visible.)}
\end{multicols}

\begin{notecode}
TICKERS = ["XLK","XLE","XLF","XLV","XLU","XLP","XLY","XLI","XLB","XLRE","XLC"]
START, END = "2004-01-01", "2024-12-31"

def load_prices(tickers, start, end):
    """Adjusted-close prices via yfinance (XLRE from 2015, XLC from 2018 — handled)."""
    import yfinance as yf
    df = yf.download(tickers, start=start, end=end, auto_adjust=True, progress=False)["Close"]
    return df[tickers]

px = load_prices(TICKERS, START, END)
rets = px.pct_change()

# 12-1 momentum signal (needs 252d history; NaN until an ETF has it)
mom = px.shift(21) / px.shift(252) - 1

print(f"Sample: {px.index[0].date()} → {px.index[-1].date()}")
print("History start per ETF:")
print(px.apply(lambda s: s.first_valid_index().date()))
\end{notecode}

% ============================================ 2 Result 1 — the IC decay curve =
\noteneedroom{150bp}
\notesection{2.\notenumsep Result 1 — the IC decay curve}

For each day, rank sectors by momentum; correlate those ranks with returns over
the \textit{next} h days. Averaging across all days gives the mean IC per
horizon — the signal's predictive power as a function of how long you hold.

\begin{notecode}
def ic_curve(mom, px, horizons, dates=None):
    """Mean rank IC (Spearman) per forward horizon."""
    out = {}
    idx = mom.index if dates is None else mom.loc[dates].index
    for h in horizons:
        fwd = px.shift(-h) / px - 1                     # forward h-day return
        ics = []
        for d in idx[::5]:                              # weekly sampling to reduce overlap
            sig, f = mom.loc[d], fwd.loc[d]
            mask = sig.notna() & f.notna()
            if mask.sum() >= 6:
                ics.append(stats.spearmanr(sig[mask], f[mask]).statistic)
        out[h] = (np.mean(ics), np.std(ics) / np.sqrt(len(ics)))
    return pd.DataFrame(out, index=["IC", "se"]).T

HORIZONS = [1, 2, 3, 5, 10, 21, 42, 63, 126]
icc = ic_curve(mom, px, HORIZONS)
print(icc.round(4))
\end{notecode}

\begin{notecode}
# fit exponential decay IC(h) = IC0 * exp(-h/tau); half-life = tau * ln 2
h = np.array(HORIZONS, dtype=float)
ic = icc["IC"].values

def expdecay(h, ic0, tau): return ic0 * np.exp(-h / tau)
(ic0, tau), _ = curve_fit(expdecay, h, ic, p0=[max(ic[0], 0.05), 30], maxfev=5000)
half_life = tau * np.log(2)

fig, ax = plt.subplots(figsize=(10, 5))
ax.errorbar(h, ic, yerr=icc["se"], fmt="o", color="steelblue", capsize=3, label="Mean rank IC")
hh = np.linspace(1, 126, 200)
ax.plot(hh, expdecay(hh, ic0, tau), "--", color="crimson",
        label=f"Exponential fit: half-life ≈ {half_life:.0f} days")
ax.axhline(0, color="black", lw=0.8)
ax.set_xlabel("Forward horizon (trading days)"); ax.set_ylabel("Information coefficient")
ax.set_title("Sector momentum: predictive power vs holding horizon")
ax.legend()
plt.tight_layout(); plt.show()

print(f"IC at 1 day:  {ic[0]:+.3f}")
print(f"IC at 63 days: {icc.loc[63,'IC']:+.3f}")
print(f"Estimated half-life: {half_life:.0f} trading days")
\end{notecode}

\textbf{Reading the curve:} the signal is worth the most immediately after
formation and gives up roughly half its power within the estimated half-life. A
strategy that rebalances slower than the half-life is trading mostly
\textit{dead} signal — this single number disciplines the turnover decision.

% ================================== 3 Result 2 — decade-over-decade erosion ==
\noteneedroom{150bp}
\notesection{3.\notenumsep Result 2 — decade-over-decade erosion}

Same signal, same universe, two eras. If momentum alpha erodes as it gets
arbitraged (crowding, publication, cheaper implementation), the recent decade
should show a flatter curve.

\begin{notecode}
era1 = mom.index[(mom.index >= "2005-01-01") & (mom.index <= "2014-12-31")]
era2 = mom.index[(mom.index >= "2015-01-01") & (mom.index <= "2024-12-31")]

icc1 = ic_curve(mom, px, HORIZONS, era1)
icc2 = ic_curve(mom, px, HORIZONS, era2)

fig, ax = plt.subplots(figsize=(10, 5))
ax.errorbar(h, icc1["IC"], yerr=icc1["se"], fmt="o-", color="steelblue", capsize=3, label="2005–2014")
ax.errorbar(h, icc2["IC"], yerr=icc2["se"], fmt="s-", color="darkorange", capsize=3, label="2015–2024")
ax.axhline(0, color="black", lw=0.8)
ax.set_xlabel("Forward horizon (trading days)"); ax.set_ylabel("Information coefficient")
ax.set_title("The same signal, two decades: alpha erosion in sector momentum")
ax.legend()
plt.tight_layout(); plt.show()

comp = pd.DataFrame({"2005–2014": icc1["IC"], "2015–2024": icc2["IC"]})
comp["change"] = comp.iloc[:, 1] - comp.iloc[:, 0]
print(comp.round(4))
\end{notecode}

% ========================================================== 4 Robustness =====
\noteneedroom{150bp}
\notesection{4.\notenumsep Robustness}

Three checks that the result isn't an artifact:

\begin{notecode}
# (a) skip-window sensitivity: 12-1 vs 12-0 vs 6-1 formation
variants = {"12-1 (base)": px.shift(21)/px.shift(252)-1,
            "12-0":        px/px.shift(252)-1,
            "6-1":         px.shift(21)/px.shift(126)-1}
rows = {}
for name, sig in variants.items():
    icc_v = ic_curve(sig, px, [5, 21, 63])
    rows[name] = icc_v["IC"]
print("IC by formation window:")
print(pd.DataFrame(rows).round(4))
\end{notecode}

\begin{notecode}
# (b) placebo: shuffle the signal in time -> IC should be ~0
mom_shuffled = mom.sample(frac=1.0, random_state=0).set_axis(mom.index)
icc_placebo = ic_curve(mom_shuffled, px, [5, 21, 63])
print("Placebo (time-shuffled signal):")
print(icc_placebo.round(4))

# (c) long-short spread: top-3 minus bottom-3 sectors, monthly rebalance
ranks = mom.rank(axis=1)
n_valid = ranks.notna().sum(axis=1)
top = ranks.ge(n_valid - 2, axis=0)          # top 3
bot = ranks.le(3, axis=0)                     # bottom 3
monthly = rets.resample("ME").apply(lambda x: (1+x).prod()-1)
sig_m = mom.resample("ME").last().shift(1)    # trade next month on last month's signal
rank_m = sig_m.rank(axis=1)
nv = rank_m.notna().sum(axis=1)
ls = (monthly.where(rank_m.ge(nv-2, axis=0)).mean(axis=1)
      - monthly.where(rank_m.le(3, axis=0)).mean(axis=1)).dropna()
for era, sl in [("2005–2014", ls.loc["2005":"2014"]), ("2015–2024", ls.loc["2015":"2024"])]:
    print(f"L/S top3-bottom3 {era}: ann. return {sl.mean()*12:+.1%}, "
          f"t-stat {sl.mean()/sl.std()*np.sqrt(len(sl)):.2f}")
\end{notecode}

% ========================================== 5 Implications & limitations =====
% Three subsections plus the references read as one closing block; at 200bp it
% split and left two bullets and the reference list alone on a 16%-full page.
\noteneedroom{330bp}
\notesection{5.\notenumsep Implications \& limitations}

\begin{multicols}{2}
\notesubsection{Implications}
\begin{notebullets}
\item \textbf{Turnover has a right answer.} With a half-life measured in weeks,
      monthly rebalancing captures most of the available signal; quarterly
      leaves much of it dead. Trading costs then decide the exact point.
\item \textbf{Capacity and expectations.} The decade-over-decade IC decline is
      what crowding looks like in data — position sizing and return
      expectations calibrated on the 2005–2014 sample would have systematically
      disappointed after 2015.
\item \textbf{Signal research must be dated.} An IC estimated over ``all
      history'' mixes two different regimes; the recent-era curve is the honest
      input for a live strategy.
\end{notebullets}

\notesubsection{Limitations}
\begin{notebullets}
\item \textbf{11 assets is a small cross-section} — rank ICs are noisy (note
      the error bars); the pattern, not any single point, is the result.
\item \textbf{Overlapping windows} — weekly sampling reduces but doesn't
      eliminate autocorrelation in IC estimates; block-bootstrap errors would
      be the rigorous upgrade.
\item \textbf{One signal, one universe} — this measures \textit{sector}
      momentum; single-stock momentum decays differently (typically slower
      formation, faster crowding).
\item \textbf{No costs} — the L/S spread is gross; at realistic ETF spreads the
      recent-era net edge thins further.
\end{notebullets}

\notesubsection{Related content}
\begin{notebullets}
\item \textbf{GameStop case study} — crowding's fast catastrophic mode; this
      note is its slow mode
\item \textbf{SMA Crossover Backtest} — turning a time-series signal into a
      disciplined strategy
\item \textbf{Walk-forward validation} (future piece) — the out-of-sample
      hygiene this note approximates with its era split
\end{notebullets}
\end{multicols}

\begin{notereferences}
\item McLean \& Pontiff (2016), ``Does Academic Research Destroy Stock Return
      Predictability?'', \textit{Journal of Finance}.
\end{notereferences}

\end{document}
